\chapter{Glossary}

\begin{description}
    \item[AS (Address Space)] A page-table context identified by an
        \texttt{AddressSpaceId} (0 = kernel, $>$0 = user). Contains the
        \texttt{TTBR0\_EL1} (user) and \texttt{TTBR1\_EL1} (kernel) pointers
        on AArch64, or \texttt{CR3} on x86-64.

    \item[Block (verb)] To remove a running thread from the scheduler and
        register its \texttt{Waker} on an \texttt{Observable}. The thread
        yields; when the event fires the waker re-admits it.

    \item[Capability] An unforgeable per-domain handle that names a kernel
        object and carries rights. Answers: ``May I perform this operation?''

    \item[Capability table] A per-AS sparse table mapping capability indices
        to kernel objects. Used for the thread table, address-space table,
        and per-domain capability table.

    \item[Completion] The terminal result of an asynchronous operation. The
        kernel tracks operation state and delivers results through polling,
        waitable completion objects, or CQ entries.

    \item[Connection capability] A client-side authority to send or call an
        endpoint. Minted from an endpoint, typically by the name service.

    \item[CQ ring (Completion Queue)] A shared-memory ring buffer for
        zero-syscall completion delivery from the kernel to userspace.

    \item[EL0] Exception Level 0 on AArch64. In CharlotteOS source and test
        names, also used as the architecture-neutral label for userspace,
        including x86-64 ring 3.

    \item[EL1] Exception Level 1 (AArch64). Kernel execution.

    \item[Endpoint] A bounded server-side IPC receive object. Has an owning
        domain, a queue capacity, and a lifecycle.

    \item[Exit-observer] An \texttt{Observer} registered on a \texttt{Thread};
        fires when the thread is dropped (exits). The ABI's intended completion
        mechanism: \texttt{submit\_worker} registers the capability as the
        worker's exit-observer.

    \item[HHDM (Higher Half Direct Mapping)] A linear mapping of all physical
        memory into the kernel's virtual address space, provided by the
        bootloader (Limine).

    \item[IdTable] A sparse \texttt{Vec<Option<T>{>}} with ID recycling. Used for
        the thread table, address-space table, and capability table.

    \item[IOMMU / SMMU] I/O Memory Management Unit. Hardware that enforces DMA
        access permissions from devices, analogous to the MMU for CPU access.

    \item[IPI (Inter-Processor Interrupt)] A signal from one LP to another.
        The kernel's IPI queues are bounded per-LP queues carrying typed
        commands.

    \item[LP (Logical Processor)] A schedulable hardware execution context,
        usually one core or hardware thread. It is distinct from a sitas shard.

    \item[Memory object] A page-backed kernel object with explicit ownership
        and capability semantics. Can be mapped, moved, lent, or copied between
        domains.

    \item[MMU] Memory Management Unit. Hardware that enforces page-table
        permissions for CPU access.

    \item[Name service] A userspace service that maps human-readable names to
        connection capabilities. The kernel does not understand service names.

    \item[Notification] ``State may have changed; inspect the corresponding
        object or queue.'' May be coalesced. Is not itself a result.

    \item[Observable / Observer] The event-notification pattern. An
        \texttt{Observable} holds \texttt{Weak<dyn Observer>} references;
        when it fires, it calls \texttt{Observer::notify()}.

    \item[PerLp<T>] A boxed slice of per-LP cells with lock wrapping.
        Interrupt-safe (masks interrupts on acquire). Use for data accessed
        from interrupt handlers or multiple LPs.

    \item[Protection domain] An address space and authority boundary. Owns a
        capability table, memory mappings, threads, and failure state. Typically
        implemented as a process.

    \item[Quantum timer] The per-LP periodic timer (10 ms default) that drives
        preemptive context switching in the \texttt{RoundRobin} scheduler.

    \item[Reaper] The deferred-thread-cleanup mechanism: dead threads are
        moved to \texttt{DEAD\_THREADS} and dropped from the next thread's
        context after the context switch.

    \item[Reply token] One-shot authority to complete a blocking or deferred
        endpoint call. Created by the kernel, bound to the caller's pending
        call, consumable exactly once, invalidated by cancellation or teardown.

    \item[Shard] A userspace ownership and execution domain. One executor runs
        its tasks; at most one task at a time mutates shard-owned state;
        cross-shard access goes through typed messages.

    \item[ShardLocal<T>] A lock-free per-LP container using
        \texttt{UnsafeCell<T>} with owner-check and borrow-flag guard.
        sitas-derived. Use for single-LP, thread-only kernel data.

    \item[ShardMailbox<M>] A typed bounded per-LP queue with
        \texttt{ShardSender<M>} (cloneable, any LP can send) and
        \texttt{ShardReceiver<M>} (single-consumer). First-class backpressure.

    \item[Supervisor] Kernel-side component that creates and destroys
        protection domains, loads ELF images, delivers bootstrap capabilities,
        and handles domain teardown.

    \item[TrapFrame] An architecture-specific snapshot of user register state at
        an exception. Syscall dispatch combines it with trusted per-thread or
        per-LP state to identify the caller.

    \item[Waker] An \texttt{Observer} wrapping a \texttt{ThreadId}. When
        notified, calls \texttt{submit\_ready\_thread(tid)}. In userspace,
        a Rust \texttt{Waker} is a hint to repoll a future --- it is not a
        kernel completion.

    \item[Yield] To voluntarily give up the LP by calling
        \texttt{yield\_lp()}. Saves the thread's context and lets the scheduler
        run the next ready thread.
\end{description}

\endinput
